\documentclass[10pt,a4paper]{amsart}
\usepackage[margin=19mm]{geometry}
\usepackage{amsmath,amssymb,mathtools}
\usepackage[scr=rsfs,cal=euler]{mathalfa}
\usepackage{xfrac}
\usepackage[unicode,hidelinks,bookmarks=false]{hyperref}
\newcommand{\ie}{i.e.\ }
\newcommand{\cf}{cf.\ }
\newcommand{\resp}{respectively\ }
\newcommand{\Subsection}{\subsection}
\providecommand{\nobreakdash}{\nobreak\hbox{-}}
\setlength{\emergencystretch}{2em}
\multlinegap0pt
\usepackage{xcolor}
\usepackage{tikz}
\usepackage{amssymb}
\usepackage{enumerate}
\usepackage[shortlabels]{enumitem}
\usepackage{float}
\newtheorem{lemma}{Lemma}
\def\E{{\mathbb E}}
\def\Ima{{\operatorname{Im}}}
\def\Var{{\operatorname{Var}}}
\newcommand{\cC}{\mathcal{C}}
\renewcommand{\phi}{\varphi}
\title{The shape of a random numerical semigroup}
\author{Maria Bras-Amor\'os, Nathan Kaplan, Deepesh Singhal}
\date{Source: arXiv:2604.26127v1 (CC BY 4.0)}
\begin{document}
\maketitle

\noindent\emph{Source and licence.} The shape of a random numerical semigroup. Version arXiv:2604.26127v1 (CC BY 4.0), distributed by arXiv under the Creative Commons Attribution 4.0 International licence, http://creativecommons.org/licenses/by/4.0/, as recorded in the arXiv licence metadata for this exact version and shown on the abstract page https://arxiv.org/abs/2604.26127v1. Full licence notice: \texttt{licenses/arxiv-2604.26127-CC-BY-4.0.txt}.\par\smallskip

\noindent\textit{Editorial scope.} Counted unit: the article's lemma A2 (source0.tex lines 1122--1135). The article's complete original proof of it is delivered with it (source0.tex lines 1136--1262, one \texttt{proof} environment of 127 lines). No mathematics is written, repaired or re-derived: the statement and the proof are the article's own lines, in the article's own notation, and no retained line is substituted or re-flowed.  Retained uncounted as the source-local context the proof needs: lines 1105--1115.  Nothing was omitted from any retained span, so the package carries no omission record.  No retained line contains a citation, so the wrapper carries no bibliography and no bibliography entry.  No retained line contains a figure environment, an image-inclusion command or drawing code, so no image is delivered and the package declares no asset dependency.  Editorial formatting only, in addition: an AMS \texttt{amsart} preamble that restates the article's own declarations and macros used by the retained lines, namely the environments \texttt{lemma} on separate counters, together with the standard packages those lines need.  The retained environments print in this document as Lemma 1 in order of appearance, whereas the article prints the same items with the numbering of its own full document.  The complete native source of the article as distributed in the arXiv e-print for this exact version is bundled unchanged as \texttt{sources/199366/source0.tex}.
\begin{lemma}\label{Lem: aks depth 3}
Consider $(A_1,A_2)\in \Ima(\chi_{g,m,l})$. For $|A_1|<k\leq 2m-g+l-|A_2|$, we have
\[
\E_g[a_{k}(S)\mid S\in \cC(g,m,l), \chi_{g,m,l}(S)=(A_1,A_2)] =m+l+\frac{(m-l)(k-|A_1|)}{2m-g+l-|A_1|-|A_2|+1},
\]
and 
\begin{align*}
&\Var_g[a_{k}(S)\mid S\in \cC(g,m,l), \chi_{g,m,l}(S)=(A_1,A_2)] \\
&=\frac{(m-l)(g-m-2l-1+|A_1|+|A_2|) (k-|A_1|) (2m-g+l-|A_2|+1-k)}{(2m-g+l-|A_1|-|A_2|+1)^2 (2m-g+l-|A_1|-|A_2|+2)}.  
\end{align*}
\end{lemma}

\begin{lemma}\label{Lem: depth 3 fix A1 A2}
Assume $\epsilon_1< \frac{1}{6\sqrt{5}}$ and $g>\frac{1}{\epsilon_1}$.
Suppose $g,m,l$ satisfy $|m-\gamma g|<\epsilon_1 g$ and $l<\epsilon_1 g$.  Also suppose that $(A_1,A_2)\in \Ima(\chi_{g,m,l})$ and $\alpha$ satisfies $3\epsilon_1< \alpha \leq \frac{1}{\sqrt{5}} -2\epsilon_1$.
Then
\[
\Bigl|\E_g \big[\psi_S(\alpha)\mid S\in \cC(g,m,l), \chi_{g,m,l}(S)=(A_1,A_2) \big]-(\gamma+\phi\alpha)\Bigr|
\leq 25\epsilon_1. 
\]
Moreover,
\[
\Bigl|\E_g \big[\psi_S(\alpha)^2\mid S\in \cC(g,m,l), \chi_{g,m,l}(S)=(A_1,A_2) \big]-(\gamma+\phi\alpha)^2\Bigr|
< 109\epsilon_1. 
\]
\end{lemma}

\begin{proof}
Let $t=1+\lfloor\alpha(g-1)\rfloor$, so $\psi_S(\alpha)= \frac{a_{t}(S)}{g}$. The condition $3\epsilon_1< \alpha \leq \frac{1}{\sqrt{5}} -2\epsilon_1$, along with $|m-\gamma g|<\epsilon_1 g$ and $l<\epsilon_1 g$ ensure that $|A_1|<t\leq 2m-g+l-|A_2|$.

Write the scaled quantities
\[
\rho=\frac{m}{g},\qquad \lambda=\frac{l}{g},\qquad
a=\frac{|A_1|}{g},\qquad b=\frac{|A_2|}{g},\qquad
\kappa=\frac{t}{g},
\]
and set
\[D=2\rho-1+\lambda-a-b+\frac{1}{g}.\]

First, since $|m-\gamma g|\leq \epsilon_1 g$, $l\leq \epsilon_1 g$, $|A_1|,|A_2|\leq l$, and $\frac{1}{g} \leq \epsilon_1$, we have
\[
|\rho-\gamma|\leq \epsilon_1,\qquad
0 < \lambda\leq \epsilon_1,\qquad
0\leq a,b\leq \lambda.
\]
Also $|\kappa-\alpha|\leq \frac{1}{g}\leq \epsilon_1$ since $t=1+\lfloor\alpha(g-1)\rfloor$.

We apply the result of Lemma~\ref{Lem: aks depth 3} not to the function $a_k(S)$ but to the function $\frac{a_k(S)}{g}$.  Dividing by this constant $g$ divides the expected value by $g$ and the variance by $g^2$.  We see that
\[
\E_g\left[\psi_S(\alpha)\mid S\in \cC(g,m,l),\ \chi_{g,m,l}(S)=(A_1,A_2)\right]
= \rho+\lambda+\frac{(\rho-\lambda)(\kappa-a)}{D}.
\]
Hence
\begin{align*}
& \Bigl| \E_g\left[\psi_S(\alpha)\mid S\in \cC(g,m,l),\ \chi_{g,m,l}(S)=(A_1,A_2)\right]-(\gamma+\phi\alpha)\Bigr| \\
\leq & |\rho-\gamma|+\lambda+\phi|\kappa-\alpha-a|
   +\Bigl|\frac{\rho-\lambda}{D}-\phi\Bigr|\cdot|\kappa-a|.
\end{align*}
We bound each of the terms in this expression.
\begin{enumerate}
    \item We know that $|\rho-\gamma|\leq \epsilon_1$ and $\lambda\leq \epsilon_1$.
    \item We know that $|\kappa-\alpha|\leq \epsilon_1$ and $a\leq \epsilon_1$, so
    \[
    \phi|\kappa-\alpha-a|\leq 2\phi\epsilon_1.
    \]
    \item We now consider the term $D$.  Since $a+b \le 2\lambda$ and $2\rho \ge 2\gamma - 2\epsilon_1$, we have
    \[
    D\geq (2\rho-1)+\lambda-2\lambda\geq (2\gamma-1)-2\epsilon_1-\epsilon_1
    = \frac1{\sqrt5}-3\epsilon_1 \ \ge\ \frac1{2\sqrt5},
    \]
    because $\epsilon_1<\frac{1}{6\sqrt5}$. Hence $D^{-1}\leq 2\sqrt5$.
    \item Next,
    \begin{align*}
    \Bigl|\frac{\rho-\lambda}{D}-\phi\Bigr|
    &=\frac{\big|(\rho-\lambda)-\phi(2\rho-1+\lambda-a-b+\tfrac1g)\big|}{D}\\
    &=\frac{\big|(1-2\phi)(\rho-\gamma)-(1+\phi)\lambda+\phi(a+b)-\tfrac{\phi}{g}\big|}{D}\\
    &\leq \frac{(2\phi-1)\epsilon_1 +(1+\phi)\epsilon_1 +2\phi\epsilon_1+\phi\epsilon_1}{D}
    = \frac{6\phi\epsilon_1}{D}.
    \end{align*}

Applying the bound $D^{-1}\leq 2\sqrt5$ from (3), we see that
    \[
    \Bigl|\frac{\rho-\lambda}{D}-\phi\Bigr|\leq 12\phi\sqrt{5}\epsilon_1.
    \]
    Moreover,
    \[
    |\kappa-a|\leq \kappa +a 
    \leq \left(\frac{1}{g}+\alpha\right) +\epsilon_1 
    \leq \left(\epsilon_1 +\frac1{\sqrt5}-2\epsilon_1\right)+\epsilon_1
    = \frac1{\sqrt5}.
    \]
    This implies
    \[    
    \Bigl|\frac{\rho-\lambda}{D}-\phi\Bigr|\cdot|\kappa-a|
    \leq 12\phi\epsilon_1.
    \]
\end{enumerate}
Collecting (1)-(4), we see that
\begin{align*}
& \Bigl| \E_g\left[\psi_S(\alpha)\mid S\in \cC(g,m,l),\ \chi_{g,m,l}(S)=(A_1,A_2)\right]-(\gamma+\phi\alpha)\Bigr| \\
 \leq & \epsilon_1+\epsilon_1+2\phi\epsilon_1+12\phi \epsilon_1
= (2+14\phi)\epsilon_1 <25\epsilon_1.
\end{align*}
This proves the claimed bound for the expectation.



Dividing the formula for the variance given in Lemma~\ref{Lem: aks depth 3} by $g^2$ gives
\begin{align*}
&\Var_g[\psi_S(\alpha)\mid S\in\cC(g,m,l),\ \chi_{g,m,l}(S)=(A_1,A_2)]\\
&=\frac{(\rho-\lambda)\bigl(1-\rho-2\lambda-\tfrac1g+a+b\bigr)
(\kappa-a)\bigl(2\rho-1+\lambda-b+\tfrac1g-\kappa\bigr)}
{gD^2(D+\tfrac1g)}.    
\end{align*}
We bound the numerator and denominator using the hypotheses.


Since $D\geq \frac{1}{2\sqrt{5}}$, we know that
\[D^2(D+\tfrac1g)\ \ge\ \Big(\frac1{2\sqrt5}\Big)^3.\]

We bound each factor of the numerator from above:
\[
\rho-\lambda\ \le\ \gamma+\epsilon_1,\qquad
1-\rho-2\lambda-\tfrac1g+a+b\ \le\ 1-\rho\ \le\ 1-\gamma+\epsilon_1,
\]
\[
\kappa-a\ \le\ \alpha+\tfrac1g\ \le\ \frac1{\sqrt5},\qquad
2\rho-1+\lambda-b+\tfrac1g-\kappa\ \le\ 2\rho-1+\lambda+\tfrac1g
\ \le\ \frac1{\sqrt5}+4\epsilon_1.
\]
Therefore
\begin{align*}
&\Var_g[\psi_S(\alpha)\mid S\in\cC(g,m,l),\ \chi_{g,m,l}(S)=(A_1,A_2)]\\
\le &\ \frac{1}{g}(2\sqrt5)^3
(\gamma+\epsilon_1)(1-\gamma+\epsilon_1)\Big(\frac1{\sqrt5}\Big)
\Big(\frac1{\sqrt5}+4\epsilon_1\Big).
\end{align*}
We evaluate the right-hand side at the largest $\epsilon_1$ allowed by our hypotheses, 
$\epsilon_1=\frac{1}{6\sqrt5}$, and see that
\[
\Var_g[\psi_S(\alpha)\mid S\in\cC(g,m,l),\ \chi_{g,m,l}(S)=(A_1,A_2)]<\frac{9}{g}< 9\epsilon_1.
\]

Now,
\begin{align*}
&\  \Bigl|\E_g \big[\psi_S(\alpha)^2\mid S\in\cC(g,m,l),\ \chi_{g,m,l}(S)=(A_1,A_2)\big]-(\gamma+\phi\alpha)^2\Bigr|\\
= &\  \Bigl| \Var_g[\psi_S(\alpha)\mid \cC(g,m,l),\ \chi_{g,m,l}] +\E_g \big[\psi_S(\alpha)\mid \cC(g,m,l),\ \chi_{g,m,l}\big]^2 -(\gamma+\phi\alpha)^2\Bigr|\\
\leq &\ \Bigl| \Var_g[\psi_S(\alpha)\mid \cC(g,m,l),\ \chi_{g,m,l}] \Bigr|\\
 &+\Bigl|\E_g \big[\psi_S(\alpha)\mid \cC(g,m,l),\ \chi_{g,m,l}\big] -(\gamma+\phi\alpha)\Bigr|
\cdot \Bigl|\E_g \big[\psi_S(\alpha)\mid \cC(g,m,l),\ \chi_{g,m,l}\big] +(\gamma+\phi\alpha)\Bigr|\\
< &\ 9\epsilon_1 + 25\epsilon_1 (2+2)
=109\epsilon_1.\qedhere
\end{align*}
\end{proof}

\end{document}
